\chapter{Introduction}

In this project, a Python programming environment is created targeted at new programmers. The environment builds upon the idea of helping users learn by removing some of the complexity behind programming, in turn creating a less overwhelming experience. In particular, this project focuses on heavily restricting the expressiveness of Python, then gradually increasing it in line with the user's knowledge. This is done by creating an environment that restricts both language features and built-in functions.

*AT THE END ADD SOMETHING ABOUT HOW WELL THE PROJECT DOES THIS*

\section{Definitions}

A \textit{restriction} consists of removing syntax from the language; code containing restrictions cannot be run, just as code containing a syntax error cannot. Restrictions come in two forms: 
\begin{itemize}
    \item \textit{Language features} or just \textit{features}: 
    Programming language constructs and data structures (e.g. loops, conditionals, functions, lists).
    \item \textit{Built-in functions} or just \textit{built-ins}: The set of Python built-in functions defined in the Python documentation\footnote{\href{https://docs.python.org/3/library/functions.html}{https://docs.python.org/3/library/functions.html}}.
\end{itemize}

\section{Motivation}

One problem many novice programmers face is how overwhelming programming can feel. The combination of learning many abstract concepts, as well as a completely new syntax for representing ideas, both lead to difficulty \cite{green_problems_1980,lahtinen_study_2005}. Creating new tools to help them overcome these challenges has been a focus of the computing education field. The idea behind many tools is to create a more rigid programming environment, in the hope that a more constrained environment will leave the learner more focused and less overwhelmed [CITE].
Students often struggle the most with understanding exactly how a computer takes code and reaches the output \cite{lahtinen_study_2005}. Notional machines help explain why; a notional machine is a simplified mental model of a computer which helps learners think through how a program is executed \cite{sorva_notional_2013}. When users are faced with a large, daunting syntax, there is a mismatch between their notional machine and the full capabilities of the language. Thus, difficulty emerges in explaining behaviour, as the user's concept of how the computer works does not match reality \cite{hermans_hedy_2020}.
Adding restrictions reduces this discrepancy.

Furthermore, restrictions can be used to help learners keep advancing their skills. By removing known features, students can be forced to understand new concepts or think about programming in creative ways. 

\section{Related Work}

Numerous programming learning environments have been created. Discussed here are two that directly implement the approaches used in the project.

Hedy\footnote{\href{https://hedy.org/}{https://hedy.org/}} builds directly on the notional machine concept above. It is a programming language, rather than just an environment, but the process of learning it takes a gradual approach- concepts are introduced one at a time. In addition, Hedy has its own multilingual syntax \cite{hermans_hedy_2020}. It uses natural language keywords in the learner's native language. This is an interesting approach to teaching coding; whilst \citet{lahtinen_study_2005} finds syntax to be one of the easier elements to learn, using familiar syntax reduces the overall cognitive load [CITE]. Consequently, learners can focus more on understanding program flow. However, this approach has drawbacks. Hedy is intended for learning; at some point, students must transition to a `real' programming language. This can come with difficulties... These problems are not unique to Hedy but are a symptom of tools that create new teaching languages, for example, Scratch\footnote{\href{https://scratch.mit.edu/}. As such, whilst the gradual approach of Hedy was adopted, it was decided to use Python rather than a new language to prevent this transition.

//////////////

Another avenue in research is the use of visual features, which have long been shown to aid learning. This approach has been explored for programming purposes: for example, visualising code execution is the main idea behind learning environments such as Python Tutor\footnote{\href{https://pythontutor.com/}{https://pythontutor.com/}} or Thonny\footnote{\href{https://thonny.org/}{https://thonny.org/}}. Both programs make use of visuals in different ways. Python Tutor lays out all the necessary information a learner should track, reducing the intermediate values they have to memorise; Thonny isolates function calls by creating separate windows to visualise each one. The efficacy of adding visual features to learning environments thus makes a strong case for their inclusion in the project.
